\section{Extended PDE Transplantation Notes}
\label{app:pde}

\subsection{General packet-basis recipe}

The semi-discrete transplantation principle becomes operational once one fixes a concrete packet basis. After Section~\ref{sec:continuous}, this recipe gives the discretization-side implementation of the continuous reduction theorem:
\begin{enumerate}
\item discretize the PDE solver on a chosen grid or mesh;
\item choose an analysis map $V_h$ and synthesis map $U_h$ for a finite packet frame;
\item compute or estimate the packet-basis matrix $\mathcal K_h=V_hK_hU_h$;
\item fit or approximate that matrix by a transported-symbol form plus residual;
\item train \MiNO\ so that its learned transport, symbol, and residual budgets track that fitted decomposition.
\end{enumerate}
This finite recipe asks for a tractable packet-basis object that reflects the microlocal geometry identified by the continuous reduction theorem.

\subsection{Helmholtz: principal symbol and packet transport}

For the variable-coefficient Helmholtz operator
\[
P(x,D)u = -\nabla\cdot(c(x)\nabla u) - \omega^2 n(x)u,
\]
the principal symbol is
\[
p(x,\xi)=c(x)|\xi|^2 - \omega^2 n(x).
\]
At high frequency one expects singularities and oscillatory energy to follow the Hamiltonian flow generated by $p$, at least over short scales before stronger interference and caustic effects dominate.  The inverse Helmholtz map, however, is a resolvent rather than a pure propagator.  Its packet-basis solver should therefore have the following structure:
\begin{enumerate}
\item one or more branchwise reindexings of packet mass along outgoing canonical relations;
\item a near-characteristic/radiation symbol envelope reflecting the local factor $(p+i0)^{-1}$ after microlocal cutoff;
\item a local amplitude correction reflecting lower-order or Jacobian contributions;
\item a residual accounting for leakage, reflection, branch splitting, and finite-frame truncation.
\end{enumerate}
This is the branch-factored transported-symbol-plus-resolvent decomposition used in the Helmholtz corollary.  Treating high-$k$ Helmholtz as only a single transported-symbol law is too optimistic; the symbol and routing budgets are load-bearing.

\paragraph{Hamiltonian picture.}
If one introduces a semiclassical parameter $h\sim \omega^{-1}$ and writes the principal symbol in semiclassical form, the associated bicharacteristic system is
\[
\dot x = \partial_\xi p(x,\xi),\qquad
\dot \xi = -\partial_x p(x,\xi).
\]
This is the source of the packet transport law. The canonical branch of \MiNO\ learns a discrete approximation to this motion instead of a generic token-mixing rule.

\paragraph{Why the residual and symbol budgets are nontrivial.}
Helmholtz is also the setting in which residual, routing, and resolvent-symbol budgets become genuinely important. Even away from trapping, one encounters reflection, turning behavior, outgoing/incoming separation, and packet leakage due to finite frame resolution. Near caustic precursors or branch splitting, a single-destination packet law may be too coarse. Near the characteristic shell, a smooth order-only multiplier may underfit the outgoing resolvent envelope. The residual term in the present theorem should therefore be read as a mathematically honest placeholder for these effects only after the explicit $E_{\mathrm{route}}$, $E_{\mathrm{cross}}$, and $E_{\mathrm{rad}}$ budgets have been accounted for.

\subsection{Wave propagation: short-time propagators}

For a short-time wave propagator, the packet-basis matrix should be even closer to the transported-symbol form than in the Helmholtz inverse setting. The reason is that propagation itself is the leading behavior. In practice, this means the wave benchmark should be the most direct place to read off the transport budget:
\begin{itemize}
\item if transport is right and symbol is roughly right, the model should preserve wavefront geometry;
\item if transport is wrong, phase error and packet consistency should reveal this almost immediately.
\end{itemize}

\paragraph{Short-time advantage.}
The short-time wave propagator is especially favorable because the canonical graph remains locally simple over short horizons. This is why the first serious empirical target for \MiNO\ should be short-time or moderate-time wave propagation rather than very long-time dynamics with repeated branch interactions.

\paragraph{Expected diagnostic behavior.}
If the transport branch is correct, the packet cloud should shift coherently and phase error should remain low. If the transport branch is wrong but the symbol branch is strong, one may still see moderate field accuracy at low frequency while the packet-consistency and local phase diagnostics degrade sharply. This is the cleanest demonstration regime for the theorem-aware metrics.

\subsection{Elliptic smoothing as identity-canonical symbol action}

Elliptic maps such as Darcy-type operators emphasize a different part of the architecture. The packet-basis representation is expected to be dominated by an identity-canonical pseudodifferential symbol with weak packet relocation. This motivates a symbol branch local in $(x,\xi)$: a purely transported architecture overfits the wrong primitive in this regime, while a purely translation-invariant spectral architecture underfits variable-coefficient elliptic symbols. The microlocal design is mixed on purpose.

\paragraph{Principal-symbol reading.}
For a uniformly elliptic symbol $p(x,\xi)$ with no hyperbolic propagation structure in the relevant regime, there is no reason to expect large packet advection across phase space. The dominant effect is symbol smoothing through a parametrix symbol $q(x,\xi)\sim p_m(x,\xi)^{-1}+\cdots$. In this regime the canonical relation is the identity, the transport map should be close to the identity, and the most meaningful theorem budget is the local symbol budget together with the smoothing residual rather than the transport budget.

\paragraph{Role of this regime.}
Darcy-like smoothing tasks test whether the packet calculus specializes to identity-canonical pseudodifferential behavior when nontrivial transport is not the leading mechanism.

\subsection{Navier--Stokes and splitting intuition}

A one-step linearization of an incompressible flow update can be understood schematically as
\[
u_{t+\Delta t}
\approx
e^{\Delta t(\nu \Delta - b(x)\cdot\nabla)}u_t
 + \text{pressure / nonlinear residual}.
\]
In packet space:
\begin{enumerate}
\item the transport part $b(x)\cdot\nabla$ moves packets;
\item the diffusion part $\nu\Delta$ damps them by dissipative symbol factors;
\item the remaining nonlinear pressure/closure effects enter the residual channel.
\end{enumerate}
This makes Navier--Stokes an especially natural mixed-regime benchmark for \MiNO.

\paragraph{Linearized symbol picture.}
The advection part contributes a Hamiltonian-like transport term, while the viscous part contributes a negative real symbol proportional to $|\xi|^2$. Thus the same update can naturally require both packet transport and frequency-dependent damping. In a variable-coefficient setting the damping should be read as a dissipative local symbol $m_t(x,\xi)\approx\exp[-tq(x,\xi)]$, not merely as a global Fourier multiplier. This is exactly the kind of mixed regime that point-space or global-frequency-only primitives tend to entangle.

\paragraph{Why one-step updates are the right first target.}
The present theorem and architecture are easiest to interpret on one-step or short-horizon updates. Full long-horizon Navier--Stokes rollouts introduce compounding nonlinear closure errors and solver instability issues that are orthogonal to the main architectural claim.

\subsection{PDE-wise budget stressors}

For completeness, we record the main theorem budgets that become active when a PDE family deviates from the ideal single-destination transported-symbol law.

\paragraph{Helmholtz.}
Trapping, turning points, branch splitting, reflection, caustic formation, and near-characteristic resolvent growth can enlarge the residual or the explicit routing/radiation budgets.

\paragraph{Wave propagation.}
Long-time evolution, boundary interaction, and multibranch interference make the packet law less sparse.

\paragraph{Darcy-type elliptic maps.}
The main issue is not transport complexity but rather whether the packet representation is unnecessarily elaborate for a smoothing-dominant map.

\paragraph{Navier--Stokes.}
Nonlinear closure, pressure projection, and aliasing between scales can all create residual terms not captured by a simple transport-plus-symbol split.

This list is important because it shows that the residual and extension terms are not generic garbage terms. They have PDE-specific mathematical content and motivate the branch, resolvent-symbol, and transport--diffusion hooks in the executable model.

\subsection{From continuous reduction to finite transplantation}

The relation between the continuous theorem of Section~\ref{sec:continuous} and the finite transplantation principle can be written as a three-step map.
\begin{enumerate}
\item prove or assume a continuous packet reduction for the target propagator;
\item discretize the propagator in a packet frame, producing $\mathcal K_h$;
\item compare the learned \MiNO\ layer to the reference transported-symbol part of $\mathcal K_h$ using Theorem~\ref{thm:finite-transported-symbol}.
\end{enumerate}
This is the exact sense in which the present paper unifies the continuous and discrete viewpoints. The continuous theorem identifies the correct operator structure; the finite theorem gives the machine-checkable error law once that structure has been discretized.

\subsection{Recommended PDE order for the project}

The mathematics and engineering should proceed in the following order.

\paragraph{First: wave propagation.}
This is the clearest transport regime and the cleanest place to validate the packet viewpoint.

\paragraph{Second: variable-coefficient Helmholtz.}
This is the most important oscillatory elliptic regime and the one with the strongest novelty upside.

\paragraph{Third: Darcy-type elliptic maps.}
These provide the necessary reduction check that \MiNO\ does not collapse on smooth problems.

\paragraph{Fourth: one-step Navier--Stokes.}
This gives the mixed advection-diffusion regime and the bridge to broader scientific-ML benchmarks.

This order is not merely pragmatic. It matches the theorem's interpretability: start where transport is cleanest, then move toward mixed and less ideal regimes.

\subsection{PDE-to-budget correspondence table}

Table~\ref{tab:pdebudget} summarizes how the theorem budgets should be read in each PDE family.

\begin{table}[t]
\centering
\caption{PDE-to-budget correspondence for the present project.}
\label{tab:pdebudget}
\begin{tabular}{p{0.21\textwidth}p{0.22\textwidth}p{0.20\textwidth}p{0.25\textwidth}}
\toprule
PDE family & Dominant budget & Main residual source & Best diagnostic \\
\midrule
Short-time wave & $\eps_{\mathrm{tr}}$ & multibranch spill, boundaries & phase error + packet consistency \\
Variable-coefficient Helmholtz & $\eps_{\mathrm{tr}}+E_{\mathrm{route}}+E_{\mathrm{rad}}$ then $\eps_{\mathrm{res}}$ & caustics, reflection, turning, resolvent shell growth & frequency-shift robustness + branch/radiation diagnostics \\
Darcy-type elliptic & identity-canonical $\eps_{\mathrm{sym}}$ + smoothing residual & overexpressive transport / weak symbol structure & packet-diagonal concentration and smooth-regime competitiveness \\
One-step Navier--Stokes & mixed $\eps_{\mathrm{tr}}$ and $\eps_{\mathrm{sym}}$ & closure and pressure residual & ablation asymmetry across transport and symbol branches \\
\bottomrule
\end{tabular}
\end{table}

This table explains the separation of transport, symbol, and residual terms. Different PDE families stress different theorem budgets, and the benchmark reflects that structure.

\subsection{Requirements for a stronger continuous theorem}

The continuous microlocal reduction theorem can be strengthened along four axes.

\paragraph{1. A packet almost-diagonalization theorem for the chosen frame.}
The manuscript assumes the standard off-graph decay estimate for the wavepacket frame. A stronger version proves that theorem directly for the exact packet system used in the implementation.

\paragraph{2. Sharper PDE-family constants.}
A sharper statement tracks constants in frequency, time horizon, and coefficient regularity.

\paragraph{3. A discrete-to-continuous stability map.}
The learned finite model would need to be compared simultaneously to the discrete solver and to the continuous operator with explicit consistency rates.

\paragraph{4. A formalization strategy beyond the current Lean scope.}
This requires a staged formalization of functional analysis, oscillatory integrals, and selected H\"ormander-calculus components.

These requirements define the mathematical extension path beyond the classical continuous bridge and machine-checked finite spine.
